\documentclass[12pt,a4paper]{article}
\usepackage[utf8]{inputenc}
\usepackage[T1]{fontenc}
\usepackage{lmodern}
\usepackage{graphicx}
\usepackage{setspace}
\usepackage{geometry}

\begin{document}

\begin{center}
    \centering

    {\huge \textbf{Análise, Projeto e Desenvolvimento Ágil}}\\[0.5cm]

    {\Huge \textbf{Questões de Scrum e Kanban}}\\[0.5cm]

    {\large Aluno 01: \textbf{Henrique Maia Cardosa}}\\[0.5cm]
    {\large Aluno 02: \textbf{João Miguel de Castro Menna}}\\[0.5cm]
    {\large Aluno 03: \textbf{Ricardo Gabriel Fialho Santos}}\\[2.5cm]
\end{center}

\pagebreak

\section*{Questões de Scrum}

\paragraph{}

1. R: Um caminho de resolução aderente ao Scrum seria o seguinte: o Scrum Master
promove uma conversa entre stakeholder, Product Owner e Time de Desenvolvimento para
esclarecer impacto e riscos. Se a funcionalidade for realmente crítica, o Product Owner pode
negociar com o Time de Desenvolvimento a possibilidade de cancelar a Sprint e iniciar uma
nova, com o escopo revisado. Caso contrário, o pedido deve ser registrado no Product
Backlog para priorização futura. Assim, mantém-se a transparência, a inspeção e a
adaptação, princípios centrais do controle empírico de processo.

Se o time ceder à pressão e aceitar incluir a funcionalidade sem respeitar o
processo, podem surgir alguns anti-padrões. O primeiro é a quebra do foco da Sprint, em
que o time perde previsibilidade e compromete as metas definidas, prejudicando a inspeção
dos resultados. O segundo é a erosão do Definition of Done, quando, para “dar conta” do
novo escopo, o time reduz a qualidade técnica ou deixa entregas incompletas,
comprometendo a confiabilidade do incremento. Esses problemas mostram como ignorar os
papéis e responsabilidades definidos no Scrum enfraquece o próprio ciclo empírico de
inspeção e adaptação, levando a decisões baseadas em pressão externa e não em
evidências.

\paragraph{}

2. R: Num cenário em que um stakeholder exige a inclusão de uma grande funcionalidade em
uma Sprint já em andamento, é importante analisar os limites de autoridade de cada papel
no Scrum.

O Product Owner tem autoridade sobre o Product Backlog e, portanto, pode decidir
sobre as prioridades e negociações com stakeholders. Contudo, ele não pode impor que o
Time de Desenvolvimento aceite novos itens dentro da Sprint em curso, pois, uma vez
iniciada, o conteúdo da Sprint só pode ser alterado pelo próprio time em comum acordo.
O Time de Desenvolvimento, por sua vez, tem autoridade para decidir quanto e como
o trabalho será realizado, respeitando critérios técnicos e o Definition of Done (DoD); ele não
pode, entretanto, simplesmente ignorar o valor de negócio trazido pelo pedido.

Já o Scrum Master atua como facilitador do processo, garantindo que o framework
seja seguido e que não haja quebra de foco ou de transparência; ele não pode tomar
decisões de negócio nem de ordem técnica, mas pode intervir para lembrar o time e o
Product Owner dos limites do Scrum.

\paragraph{}

3. R: No caso de planejamento sob incerteza, o time precisa assumir que não vai conseguir
prever tudo, então o Objetivo da Sprint deve ser claro e mensurável, mas sem depender de
“entregar tudo”. Um exemplo seria: “Validar a integração inicial com o sistema externo e
garantir pelo menos um fluxo completo conforme o DoD”. Assim, o valor está na validação e
no aprendizado, não na quantidade.

Na escolha dos itens do Product Backlog, é importante equilibrar valor, risco e qualidade: dar
prioridade ao que gera mais impacto, mas também trazer cedo os itens de maior risco
técnico para reduzir incertezas. O DoD continua sendo a linha de corte, para que só se
aceite aquilo que pode ser entregue de forma estável. Para fatiar histórias, o time pode
dividir grandes funcionalidades em fluxos menores de uso, separar o essencial do opcional e
buscar entregar algo inspecionável logo de início.

Como existe dependência externa, é fundamental ter um plano B: preparar histórias
alternativas que não dependam dessa integração, de forma que, se houver atraso, o time
ainda consiga entregar valor. O critério de sucesso da Sprint não deve ser “concluir todos os
itens”, mas sim alcançar o Objetivo definido — por exemplo, permitir que stakeholders
testem pelo menos um fluxo crítico, funcionando de ponta a ponta e de acordo com o DoD.
Isso mantém a transparência, dá condições para inspeção e abre espaço para adaptação na
Sprint seguinte.

\paragraph{}

4. R: O Definition of Done (DoD) é um acordo do time que define o que significa uma entrega
estar realmente concluída, valendo para todas as histórias e garantindo qualidade
padronizada; já os critérios de aceitação são específicos de cada item do backlog e
descrevem o comportamento esperado para validar se a funcionalidade atende ao que foi
pedido. No nível de produto, o DoD assegura consistência da solução como um todo,
enquanto no nível de time ele orienta práticas técnicas, testes e revisões.
Para a história “Login básico com e-mail e senha”, um DoD mínimo poderia incluir: código
revisado e versionado, testes automatizados passando, cobertura mínima definida,
documentação atualizada, integração contínua sem falhas, usabilidade básica validada e
ausência de vulnerabilidades conhecidas. Já três critérios de aceitação testáveis seriam:
Dado um usuário cadastrado, quando insere e-mail e senha corretos, então o acesso é
concedido.

Dado um usuário cadastrado, quando insere senha incorreta, então o sistema exibe
mensagem de erro sem revelar detalhes sensíveis.
Dado um usuário não registrado, quando tenta acessar, então o sistema bloqueia a entrada
e sugere cadastro.

Mesmo que os critérios de aceitação estejam incompletos ou mal escritos, o DoD funciona
como rede de proteção, pois exige padrões técnicos (como testes, segurança,
documentação e revisão), impedindo que o incremento seja liberado sem qualidade mínima
e mantendo a confiabilidade do produto.

\paragraph{}

5. R: Um time que faz Dailies longas e sem foco, Review com incremento incompleto e Retro
irregular provavelmente sofre de três causas-raiz: falta de clareza de objetivo da Sprint,
ausência de disciplina no cumprimento do DoD e baixo comprometimento com as
cerimônias. O experimento para as próximas duas Sprints pode ser: limitar a Daily a 15
minutos com pauta fixa, reforçar na Planning e na Review que apenas incrementos prontos
(segundo o DoD) serão mostrados, e tornar a Retro obrigatória com tempo reservado no
calendário. As métricas a acompanhar seriam: cumprimento do Objetivo da Sprint,
percentual de histórias que atendem ao DoD, defeitos encontrados após a Sprint e tempo
médio de aprovação de PRs. A regra de decisão seria: se houver melhora consistente em
pelo menos dois indicadores-chave (por exemplo, mais histórias entregues dentro do DoD e
incremento estável na Review), manter a mudança; caso contrário, reavaliar e ajustar o
processo

\pagebreak

\section*{Questões de Kanban}

\paragraph{}

1. R: Os princípios formam um ciclo: tornar o fluxo visível expõe gargalos; limites de
WIP reduzem paralelismo e filas; gerenciar o fluxo com métricas permite ajustes; políticas
explícitas orientam a movimentação dos itens; cadências de feedback consolidam melhorias.
Exemplo: ao limitar WIP em “Em Progresso” de 8 para 4, o time reduz multitarefa, conclui
itens mais rápido e o lead time mediano cai (menos espera).

\paragraph{}

2. R: Pela Lei de Little: Lead time ~= WIP Throughput = 30 12 ~= 2,5 semanas/ item.
Estratégias:
(1) Reduzir WIP (limites mais baixos por coluna e política de “finalizar antes
de começar”);
(2) Aumentar throughput de qualidade removendo gargalos (ex.: pair em
Estágios críticos, automação de testes, reduzir handoffs)

\paragraph{}

3. R: O alargamento em Review/Test indica gargalo.

Ações:
\begin{itemize}
    \item Operacionais: realocar pessoas, automatizar testes.
    \item Políticas: limitar WIP nessa coluna, definir critérios de entrada claros e priorizar conclusão de bloqueados antes de iniciar novos.
\end{itemize}

\paragraph{}

4. R: Comunicação: “50\% das entregas levam até 6 dias (P50) e 85\% até 10 dias (P85).”
\begin{itemize}
    \item Para stakeholders: usar P85 como prazo seguro.
    \item Internamente: usar P50 para planejar capacidade.
    \item Itens críticos: aplicar buffers ou classes de serviço especiais.
\end{itemize}

5. R: Política: puxar primeiro bloqueados/antigos (aging alto); depois Fixed Date próximos;
reservar lane/WIP específico para Expedite (uso raro e criterioso).

Riscos do abuso de Expedite: desorganiza fluxo, aumenta variabilidade e lead time dos
demais.

Mitigação: quota limitada, aprovação formal e revisão periódica do uso.

\end{document}

